\documentclass[11pt]{article}

\usepackage[margin=1in]{geometry}
\usepackage{amsmath,amssymb,amsthm}
\usepackage{booktabs}
\usepackage{array}
\usepackage[table]{xcolor}
\usepackage{siunitx}
\usepackage{enumitem}
\usepackage[expansion=false,protrusion=true]{microtype}
\usepackage[hidelinks]{hyperref}
\usepackage{titlesec}
\usepackage{caption}

% --- F1R3FLY brand accent ---
\definecolor{f1r3red}{HTML}{D7263D}
\definecolor{f1r3ink}{HTML}{1A1A1A}
\definecolor{f1r3ash}{HTML}{6B6B6B}
\definecolor{coolrow}{HTML}{EAF5EC}
\definecolor{hotrow}{HTML}{FBE9EC}
\definecolor{headrow}{HTML}{F4E1D2}

\titleformat{\section}{\color{f1r3red}\large\bfseries}{\thesection}{0.6em}{}
\titleformat{\subsection}{\color{f1r3ink}\normalsize\bfseries}{\thesubsection}{0.6em}{}

\hypersetup{colorlinks=true,linkcolor=f1r3red,citecolor=f1r3red,urlcolor=f1r3red}

\newcommand{\near}{\operatorname{near}}
\sisetup{group-separator={,},group-minimum-digits=4}

\newtheorem{remark}{Remark}
\newtheorem{observation}{Observation}

\title{\color{f1r3red}\bfseries Sector Economics and Go-to-Market Viability\\[2pt]
\large\color{f1r3ink}\mdseries When a F1R3FLY shard pays for itself, and which markets get it there\\[2pt]
\normalsize\color{f1r3ash}\mdseries A companion to ``Rent and Shard Splitting''}
\author{\normalsize F1R3FLY Industries --- working note}
\date{\normalsize \today}

\begin{document}
\maketitle

\noindent\textit{A companion note. It reuses the cost model, the dense-packing result, and the rent lifecycle of ``Rent and Shard Splitting''~\cite{rentnote} and asks the next question: per market sector, how much monthly traffic does a shard need to be economically viable, which sectors produce it, and in what order should we deploy. All dollar figures are engineering estimates on stated assumptions; every fee is a lever, flagged as such, so the tables re-derive when real schedules replace ours.}

\section{The viability condition}
\label{sec:viability}

A shard is viable when monthly phlogiston revenue covers its cost. From the rent note, one ten-node shard costs $C \approx \$4{,}160$/month all-in, and dense elastic packing~\cite{fkk2025} lets it carry on the order of $\num{30000000}$ fee-bearing transactions per month before the split threshold binds. If a sector's transactions bear a blended fee $r_{\text{tx}}$ (precharge, activation, and rent combined, in dollars), then break-even monthly traffic is simply
\[
  T^\star \;=\; \frac{C_{\text{eff}}}{r_{\text{tx}}},
  \qquad
  C_{\text{eff}} \;=\; C \times (\text{compliance multiplier}).
\]
Risk pushes on \emph{both} sides of this ratio. A riskier sector can charge a higher $r_{\text{tx}}$, but it also raises $C_{\text{eff}}$: replication, audit, and dedicated infrastructure break the ``all co-resident on one cloud'' assumption. The useful dimensionless quantity is therefore not $T^\star$ but the \emph{break-even load}, $T^\star$ divided by per-shard capacity---the fraction of a shard a sector must fill to pay for it.

\begin{observation}[Capacity is never the gate; demand is.]
Every sector below is feasible on a single shard's \emph{capacity}---even the highest-volume case sits far under the split threshold per shard. Viability is gated entirely by whether a sector produces enough fee-bearing transactions, and the break-even load varies by more than two orders of magnitude across sectors. Riskier data is \emph{easier} to make pay, not harder.
\end{observation}

\paragraph{Common parameters.} Unless noted: base cost $C=\$4{,}160$/shard/month; per-payload network cost at dense packing $C/\num{30000000}\approx\$0.00014$; low-risk shards run dense (${\sim}\num{28000000}$--$\num{30000000}$ tx/mo), mission-critical shards run with headroom (${\sim}\num{20000000}$ tx/mo). Compliance multipliers: AI $1.2\times$, music $1\times$, biotech $1.5\times$, fintech $2.5\times$, medical $4\times$.

\section{AI (ASI Alliance) --- the growth engine}
\label{sec:ai}

Assumptions: $r_{\text{tx}}=\$0.01$ (inference calls, tool invocations, and agent-to-agent settlement all bear real value); $C_{\text{eff}}=\$4{,}992$; ${\sim}\num{10000}$ tx per agent per month at machine pace; dense packing.

\textbf{Break-even: ${\sim}\num{500000}$ tx/month ${=}$ ${\sim}1.8\%$ of a shard.} An AI shard pays for itself at two percent utilisation.

\begin{table}[h]
\centering
\caption*{\textbf{Table AI.} Agent population ramp at $10\times$/year (unbounded).}
\begin{tabular}{@{}rrrrrrr@{}}
\toprule
Year & Agents & Monthly tx & Shards & Cost & Revenue & Margin \\
\midrule
1 & \num{100}     & 1.0M  & 1   & \$4{,}992    & \$10{,}000  & 2.0$\times$ \\
2 & \num{1000}    & 10M   & 1   & \$4{,}992    & \$100{,}000 & 20$\times$ \\
3 & \num{10000}   & 100M  & 4   & \$19{,}968   & \$1.0M      & 50$\times$ \\
4 & \num{100000}  & 1.0B  & 36  & \$179{,}712  & \$10M       & 56$\times$ \\
5 & \num{1000000} & 10B   & 358 & \$1.79M      & \$100M      & 56$\times$ \\
\bottomrule
\end{tabular}
\end{table}

The defining trait is the absence of a TAM ceiling. Machine-generated volume is unbounded and sticky; margin saturates near $56\times$ and then revenue simply scales. AI both load-tests the dispatch/\textsc{comm} hot path under real volume \emph{and} is the only sector whose growth does not terminate---which is why it should anchor the rollout rather than a low-fee consumer sector.

\section{Fintech (Boring Financial) --- the density engine}
\label{sec:fintech}

Assumptions: $r_{\text{tx}}=\$0.25$ blended (settlement, contract-lifecycle, and loan-servicing events); $C_{\text{eff}}=\$10{,}400$ (custody, audit, regulatory load); safety headroom.

\textbf{Break-even: ${\sim}\num{42000}$ tx/month ${=}$ ${\sim}0.2\%$ of a shard}---the most revenue-dense sector; a fintech shard pays for itself at one-fifth of one percent utilisation.

\begin{table}[h]
\centering
\caption*{\textbf{Table FIN.} Transaction ramp.}
\begin{tabular}{@{}rrrrrr@{}}
\toprule
Year & Monthly tx & Shards & Cost & Revenue & Margin \\
\midrule
1 & 50k   & 1  & \$10{,}400  & \$12{,}500  & 1.2$\times$ \\
2 & 500k  & 1  & \$10{,}400  & \$125{,}000 & 12$\times$ \\
3 & 5M    & 1  & \$10{,}400  & \$1.25M     & 120$\times$ \\
4 & 50M   & 3  & \$31{,}200  & \$12.5M     & 400$\times$ \\
5 & 500M  & 25 & \$260{,}000 & \$125M      & 481$\times$ \\
\bottomrule
\end{tabular}
\end{table}

\begin{remark}[The most assumption-sensitive number in the note.]
$\$0.25$/tx is right for high-value settlement events, where it is trivial against notional, and wrong by one to two orders of magnitude for high-frequency micropayments, where fintech behaves like AI on margins. The real schedule is basis-points-on-notional for large events and sub-cent for small ones, so fintech is bimodal; model the two instrument classes separately once Boring supplies their mix.
\end{remark}

\section{Medical (VA) --- the cohort-capped annuity}
\label{sec:medical}

The Veterans Health Administration serves roughly \num{9100000} enrolled veterans~\cite{vha2026}; the $80\%$ target is ${\sim}\num{7280000}$ patients. Assumptions: $\tau=50$ provenance-logged record accesses per patient per month; $R=\$1$/patient/month in platform fees (a conservative $2$--$7\%$ of full-EHR per-patient spend); $C_{\text{eff}}=\$16{,}640$; headroom.

\textbf{Break-even: ${\sim}\num{16640}$ patients against a ${\sim}\num{400000}$-patient shard ${=}$ ${\sim}4\%$ load.} A medical shard pays for itself at four percent utilisation.

\begin{table}[h]
\centering
\caption*{\textbf{Table MED.} $10\times$/year rollout against a finite cohort (saturates at year~4).}
\begin{tabular}{@{}rrrrrrr@{}}
\toprule
Year & Patients & Monthly tx & Shards & Cost & Revenue & Margin \\
\midrule
1 & \num{7280}     & 0.36M & 1  & \$16{,}640  & \$7{,}280     & 0.4$\times$ (loss) \\
2 & \num{72800}    & 3.6M  & 1  & \$16{,}640  & \$72{,}800    & 4.4$\times$ \\
3 & \num{728000}   & 36M   & 2  & \$33{,}280  & \$728{,}000   & 21.9$\times$ \\
4 & \num{7280000}  & 364M  & 19 & \$316{,}160 & \$7{,}280{,}000 & 23.0$\times$ \\
\rowcolor{coolrow}5--10 & \num{7280000} & 364M & 19 & \$316{,}160 & \$7{,}280{,}000 & 23.0$\times$ (flat) \\
\bottomrule
\end{tabular}
\end{table}

Two facts dominate. First, the cohort is a hard wall: a literal $10\times$/year ramp saturates the finite population by year~4, so medical revenue is \emph{logistic}---explosive then flat---and past saturation it is an annuity (${\sim}\$87$M/year revenue on ${\sim}\$3.8$M/year infrastructure), not a growth story. Cumulative decade revenue is ${\sim}\$620$M, of which infrastructure is a ${\sim}4\%$ rounding error. Second, the constraint is never economic: even tripling cost for high-availability replication ($3\times\!\times4\times$ loading) leaves ${\sim}7$--$8\times$ margin at cohort. The binding constraints are regulatory and operational.

\begin{remark}[Risk sets the rent posture.]
``\textsc{comm}s delivered while funded'' (the rent note's mortality semantics) is a safe default for music but a safety failure for a patient-record access that silently never fires on fund depletion. Mission-critical residuals must be pinned to committed on-chain balance (germ line), never allowed to deplete to eviction. This is the strongest reason medical deploys last: the rent model itself is most stressed exactly where the data is least forgiving.
\end{remark}

\section{Music --- the cold-tail fill}
\label{sec:music}

A stable of ten independent artists at ${\sim}1$M non-overlapping followers, each audient streaming a conservative ten songs/month, is $100$M plays/month. Per-play on-chain, that is ${\sim}\num{4}$ shards ($\$16{,}640$/mo), and break-even is
\[
  r_{\text{tx}}^\star \;=\; \frac{\$16{,}640}{100\text{M}} \;\approx\; \$0.000166\ \text{per play}.
\]
Any network fee above ${\sim}1/60$ of a cent per play is profitable; at $\$0.0003$/play the stable runs ${\sim}1.8\times$.

\begin{table}[h]
\centering
\caption*{\textbf{Table MUS.} Margin on the ten-artist stable (100M plays/mo, 4 shards, \$16{,}640/mo).}
\begin{tabular}{@{}rrrr@{}}
\toprule
Network fee/play & Monthly revenue & Cost & Margin \\
\midrule
\$0.0001  & \$10{,}000  & \$16{,}640 & 0.6$\times$ (underwater) \\
\$0.000166 & \$16{,}640 & \$16{,}640 & break-even \\
\$0.0003  & \$30{,}000  & \$16{,}640 & 1.8$\times$ \\
\$0.001   & \$100{,}000 & \$16{,}640 & 6.0$\times$ \\
\bottomrule
\end{tabular}
\end{table}

Two design facts. Batching plays off-chain and settling (say $100{:}1$) collapses $100$M plays to $1$M transactions on one shard at ${\sim}7\times$ margin---strictly better economics \emph{unless per-play provenance is the product}. If transparent per-play royalty and capability-attenuated rights are the differentiator over incumbents, we \emph{want} the $100$M on-chain transactions, and music becomes the abundant \emph{cold-tail packing} that monetises the spare leaves of the AI and fintech shards (rent note, bid-rent stratification). Either way music is viable; the choice is provenance granularity versus margin. The caveat is concentration: $100$M/month resting on ten artists is roster-dependent and churns, which is why music is high-margin fill, never the anchor.

\section{Web 2.0 --- participatory advertising}
\label{sec:web2}

Free-at-point-of-use consumer properties have no per-transaction user fee, so the framework above appears to fail---until advertising is made \emph{participatory}. End users pay in tokenised impressions per data payload until their usage is paid off, then either switch to impression-free payloads or accumulate advertising tokens spendable in other markets. This both closes the monetisation gap and inverts the extraction relationship.

The mechanism rests on one ratio. A data payload costs $\$0.00014$ at dense packing; an ordinary display impression at a ${\sim}\$3$ CPM~\cite{gdn2026} is worth $\$0.003$:
\[
  \boxed{\ \text{one ad impression covers the network cost of } {\sim}20 \text{ data payloads.}\ }
\]
Because impressions are tokenised on-chain, each impression is itself a fee-bearing transaction that funds the payload it rides plus margin: the ad layer is self-funding volume, not overhead. When cumulative impressions exceed cumulative usage, the user crosses into the impression-free or net-earner fork.

\begin{table}[h]
\centering
\caption*{\textbf{Table WEB.} Three Google properties; payloads as emails/sessions/edits.
Volumes from Gmail~\cite{gmail2026}, Maps~\cite{maps2026}, Workspace~\cite{workspace2026}.}
\begin{tabular}{@{}lrrrr@{}}
\toprule
Property & Payloads/mo & Shards & Infra cost/mo & Break-even ads/user/day \\
\midrule
Gmail (\num{1800000000} users) & ${\sim}3.6$T & ${\sim}\num{120000}$ & ${\sim}\$500$M & ${\sim}3$ \\
Maps (\num{2000000000} users)  & ${\sim}3.0$T & ${\sim}\num{100000}$ & ${\sim}\$416$M & ${\sim}2$--$3$ \\
Docs (\num{3000000000} users)  & ${\sim}600$B & ${\sim}\num{20000}$  & ${\sim}\$83$M  & ${\sim}0.3$ \\
\bottomrule
\end{tabular}
\end{table}

Every property is funded if a user views two or three ads per day---far below real exposure---so the typical user crosses into accumulating spendable tokens. \textbf{That is the inversion}: in incumbent Web 2.0 the user is the product, their attention extracted and the surplus captured by the platform; here the user covers their own cost with attention and \emph{keeps the surplus} as portable tokens. The network takes a thin per-payload margin at astronomical volume---a $\$0.00002$/payload take across Gmail's $3.6$T payloads is still ${\sim}\$72$M/month. This is ``the community is the moat'' and the anti-enshittification thesis rendered as an income statement.

\begin{observation}[The cross-market spend saves Docs.]
Docs has the lowest break-even (${\sim}0.3$ ads/day) but the highest ad resistance---nobody wants ads in documents. The fork resolves it: users earn tokens on high-CPM, ad-tolerant surfaces (Maps' local and sponsored context) and spend them on impression-free Docs. The cross-subsidy is carried in the \emph{user's} portable tokens, structurally unlike incumbent bundling, where the cross-subsidy is the platform's lever over the user.
\end{observation}

\paragraph{Enter through a privacy wedge, not at Gmail scale.} $\num{120000}$ shards and $\$500$M/month is a hyperscaler buildout; per-payload viability is trivial but absolute capital is enormous. The wedge is a privacy-premium slice: $1\%$ of Gmail (${\sim}\num{18000000}$ users who actively want data ownership) is ${\sim}36$B payloads/month ${\approx}\num{1200}$ shards ${\approx}\$5$M/month infrastructure, trivially funded---and ``own your data, earn from your attention'' self-selects exactly those adopters. Web 2.0 is a beachhead-then-scale play entered once the participatory rails are proven on the earlier sectors.

\section{Synthesis and rollout}
\label{sec:synthesis}

\begin{table}[h]
\centering
\caption*{\textbf{Table SUM.} The five sectors plus Web 2.0.}
\setlength{\tabcolsep}{5pt}
\renewcommand{\arraystretch}{1.2}
\begin{tabular}{@{}llllll@{}}
\toprule
\rowcolor{headrow}
Sector & Revenue density & Risk & Boundedness & Break-even load & Role \\
\midrule
Web 2.0  & ad-funded (\$0.00014/pl) & low--mod & astronomical & ${\sim}3$ ads/user/day & mass-market moat (last) \\
Music    & \$0.0003/tx              & low      & unbounded, churny & ${\sim}$near full & cold-tail fill \\
AI       & \$0.01/tx                & low--mod & unbounded, sticky & ${\sim}2\%$ & \textbf{anchor / growth} \\
Fintech  & ${\sim}\$0.25$/tx        & high     & semi-bounded & ${\sim}0.2\%$ & margin engine \\
Medical  & high/patient             & highest  & hard cohort cap & ${\sim}4\%$ & cash-cow annuity \\
\bottomrule
\end{tabular}
\end{table}

The pattern that falls out: \textbf{revenue density and risk rise together, and both run opposite to volume.} The original safe-to-critical \emph{risk} ordering is correct and unchanged. Layered on top, the \emph{economic} ordering is:

\begin{enumerate}[leftmargin=1.6em,itemsep=2pt]
  \item \textbf{AI --- anchor.} Simultaneous load-testing and viability; hardens the hot path; unbounded growth.
  \item \textbf{Fintech --- margin.} Highest revenue density rewards even modest volume; builds high-assurance discipline on financially recoverable failures.
  \item \textbf{Medical --- annuity.} Highest margin per unit, capped at the cohort, deployed last so the assurance and germ-line-pinning machinery is battle-tested first.
  \item \textbf{Music --- continuous fill.} Not a phase but the cold-tail packing that monetises spare capacity across the above.
  \item \textbf{Web 2.0 --- mass-market substrate.} Entered through a privacy wedge once the participatory ad rails are proven.
\end{enumerate}

\section{Sensitivity and open levers}
\label{sec:sensitivity}

Every figure scales transparently. The fees $r_{\text{tx}}$ are the dominant levers and should be replaced with real schedules from ASI and Boring as soon as available; the fintech number in particular is bimodal and warrants two instrument classes. The per-payload cost falls linearly with packing density, so the entire Web 2.0 case strengthens with the elastic-hashing gains the rent note quantifies---for low-fee consumer volume, dense packing is the difference between profit and loss, not a refinement. The compliance multipliers are coarse stand-ins for replication, audit, and dedicated-infrastructure costs and should be replaced with per-sector infrastructure models before any commitment. Finally, the single largest non-modelled risk is customer concentration---ten artists in music, one federal buyer in medical---which the portfolio sequencing above is in part designed to diversify.

\begin{thebibliography}{9}

\bibitem{rentnote}
F1R3FLY Industries.
\newblock Rent and shard splitting: a storage-economics companion to occupancy-driven sharding.
\newblock Working note, 2026.

\bibitem{fkk2025}
M.~Farach-Colton, A.~Krapivin, and W.~Kuszmaul.
\newblock Optimal bounds for open addressing without reordering.
\newblock arXiv:2501.02305, 2025.

\bibitem{rho2005}
L.~G. Meredith and M.~Radestock.
\newblock A reflective higher-order calculus.
\newblock \emph{Electronic Notes in Theoretical Computer Science}, 141(5):49--67, 2005.

\bibitem{oci2026}
Oracle Corporation.
\newblock Oracle Cloud Infrastructure pricing (compute, block volume, and networking price lists).
\newblock \url{https://www.oracle.com/cloud/price-list/}, accessed June 2026.

\bibitem{vha2026}
U.S. Department of Veterans Affairs, Veterans Health Administration.
\newblock About VHA: enrolled-veteran population and budget.
\newblock \url{https://www.va.gov/health/aboutvha.asp}, accessed June 2026.

\bibitem{gmail2026}
Gmail usage statistics: active users and daily message volume, 2026.
\newblock Aggregated industry reporting (Statista, Litmus, Radicati), accessed June 2026.

\bibitem{maps2026}
Google Maps usage statistics: monthly active users, 2026.
\newblock Aggregated industry reporting, accessed June 2026.

\bibitem{workspace2026}
Google Workspace active-user disclosures, 2026.
\newblock Google I/O and aggregated industry reporting, accessed June 2026.

\bibitem{gdn2026}
Display advertising benchmarks 2026: Google Display Network average CPM ${\approx}\$3.12$.
\newblock Aggregated programmatic-benchmark reporting, accessed June 2026.

\end{thebibliography}

\end{document}
